\documentclass[aspectratio=169]{beamer}

% ---------- ENCODING & FONTS ----------
\usepackage[utf8]{inputenc}
\usepackage[T1]{fontenc}

% ---------- THEME & LOOK ----------
\usetheme{Madrid}
\usecolortheme{seahorse}
\usefonttheme{professionalfonts}
\setbeamertemplate{navigation symbols}{}
\setbeamercovered{transparent}
\setbeamertemplate{itemize items}[circle]
\setbeamertemplate{enumerate items}[default]
\setlength{\parskip}{0.25em}

% ---------- COLORS ----------
\usepackage{xcolor}
\definecolor{accent}{RGB}{0,92,184}
\definecolor{soft}{RGB}{233,244,255}
\setbeamercolor{structure}{fg=accent}
\setbeamercolor{title}{fg=white,bg=accent}
\setbeamercolor{frametitle}{fg=accent}
\setbeamercolor{block title}{bg=accent!12,fg=accent}
\setbeamercolor{block body}{bg=soft}

% ---------- PACKAGES ----------
\usepackage{booktabs}
\usepackage{amsmath, amssymb}
\usepackage{array}
\usepackage{multirow}
\usepackage{hyperref}
\hypersetup{colorlinks=true,linkcolor=accent,urlcolor=accent}

% ---------- SAFE FALLBACK FOR TRANSITIONS ----------
\makeatletter
\@ifundefined{transfade}{\newcommand{\transfade}[1][]{}}{}
\makeatother

% ---------- PLACEHOLDER BOX (compile-safe; no images required) ----------
\newcommand{\placeholderbox}[2][0.82\linewidth]{%
  \begin{center}
    \fbox{\begin{minipage}[c][0.46\textheight][c]{#1}
      \centering \vspace{0.25em}
      \textbf{#2}\\[0.6em]
      \emph{Insert figure/diagram here later.}
    \end{minipage}}
  \end{center}
}

% ---------- TITLE ----------
\title{\textbf{Labour Demand, Nonwage Benefits, and Quasi-Fixed Costs \vspace{.5cm}\\  }}
\subtitle{ {\large ECON 300 - Labour Economics}}
%\institute{UNBC}
\author{Muhebullah Karimzada}

\date{Winter 2026}


\begin{document}

% ============================================================
% TITLE
% ============================================================
\begin{frame}[plain]
  \titlepage
\end{frame}

% ============================================================
\section{Learning Objectives \& Main Questions}
% ============================================================

\begin{frame}{Learning Objectives}
\transfade
\begin{itemize}[<+->]
  \item Outline quasi-fixed costs of labour and how they differ from variable (per-hour) costs.
  \item Distinguish workers vs.\ hours as separate margins of labour input.
  \item Explain the role of nonwage benefits in total compensation and in labour-demand modelling.
  \item Show how quasi-fixed costs create a wedge between VMP and the wage, shaping hoarding and dynamics.
\end{itemize}
\end{frame}

\begin{frame}{Main Questions}
\transfade
\begin{itemize}[<+->]
  \item Why can payroll-tax \textbf{ceilings} reduce employment?
  \item Why are less-expensive, unskilled workers more likely to be laid off in downturns?
  \item Can limits on weekly hours or overtime \textbf{create} jobs for others?
  \item Did short-time work/worksharing help save jobs in 2008–2009?
\end{itemize}
\end{frame}

% ============================================================
\section{Nonwage Benefits \& Total Compensation}
% ============================================================

\begin{frame}{Figure 6.1: Components of Total Compensation — \textit{Placeholder}}
\transfade
\placeholderbox{Figure 6.1 — Components of total compensation.}
\end{frame}

\begin{frame}{Figure 6.2: Total Compensation, Wage, and Nonwage Benefits — \textit{Placeholder}}
\transfade
\placeholderbox{Figure 6.2 — Wage vs.\ nonwage benefits over time/levels.}
\end{frame}

\begin{frame}{Nonwage Benefits and Total Compensation}
\transfade
\begin{block}{Unconditional Cash Transfers vs.\ Transfers-In-Kind}
\end{block}
\begin{columns}[T,onlytextwidth]
\column{0.52\textwidth}
\textbf{Employee Benefits}
\begin{itemize}[<+->]
  \item Taxes
  \item Group purchase
  \item Adverse selection
\end{itemize}

\column{0.48\textwidth}
\textbf{Employer Benefits}
\begin{itemize}[<+->]
  \item Production planning
  \item Deferred compensation
\end{itemize}
\end{columns}
\end{frame}

% ============================================================
\section{Quasi-Fixed Labour Costs}
% ============================================================

\begin{frame}{Definition and Sources}
\transfade
\begin{itemize}[<+->]
  \item \textbf{Independent of hours} worked; incurred per employee.
  \item Arise from:
    \begin{itemize}[<+->]
      \item Hiring costs
      \item Training costs
      \item Dismissal costs
      \item Nonwage benefits
    \end{itemize}
\end{itemize}
\end{frame}

\begin{frame}{Variable vs.\ Quasi-Fixed Costs}
\transfade
\begin{itemize}[<+->]
  \item \textbf{Variable} labour costs: vary with hours.
  \item \textbf{Quasi-fixed:} incurred per employee, \emph{independent} of hours.
\end{itemize}
\end{frame}

\begin{frame}{Recurring vs.\ Nonrecurring}
\transfade
\begin{itemize}[<+->]
  \item \textbf{Recurring}: payroll taxes.
  \item \textbf{Nonrecurring}: hiring/orienting new employees; dismissals.
\end{itemize}
\end{frame}

\begin{frame}{General Effect of Quasi-Fixed Costs}
\transfade
\begin{itemize}[<+->]
  \item Raise MC of hiring another \textbf{worker} relative to adding \textbf{hours} on existing workers.
  \item Discourage expansion along the \textbf{employment} margin; encourage more hours per worker.
  \item Hiring continues until present value of expected future revenues equals present value of additional costs.
\end{itemize}
\end{frame}

\begin{frame}{Profit-Maximizing Employment Rule}
\transfade
\small
\[
\underbrace{H+T}_{\text{hiring+training}}
\;+\;
\sum_{t=1}^{N}\frac{W_t}{(1+R)^{\,t}}
\;\;\le\;\;
\sum_{t=1}^{N}\frac{\text{VMP}_t}{(1+R)^{\,t}}
\]
\smallskip
\textbf{Where:}\; \(H+T\) = hiring and training costs;\; \(W_t\) = wage in period \(t\);\;
\(R\) = discount rate;\; \(N\) = expected length of employment;\;
\(\text{MRP}_t (\text{VMP}_t)\) = expected value of marginal product.
\end{frame}

\begin{frame}{WORKED EXAMPLE: Quasi-fixed Costs and Training}
\transfade
\scriptsize
\begin{tabular}{lccccccc}
\toprule
\multicolumn{1}{c}{\textbf{Costs}} &
\textbf{Employment Duration (N)} &
\textbf{Fixed (H + T)} &
\textbf{Variable ($\sum W_t$)} &
\textbf{Total} &
\textbf{Benefits ($\sum \mathrm{VMP}_t$)} &
\textbf{Total Profit} &
\textbf{Yearly Profit} \\
\midrule
 & 1 year & \$15{,}000 & \$25{,}000 & \$40{,}000 & \$30{,}000 & $-{\$}10{,}000$ & $-{\$}10{,}000$ \\
 & 2 years & \$15{,}000 & \$50{,}000 & \$65{,}000 & \$60{,}000 & $-{\$}5{,}000$ & $-{\$}2{,}500$ \\
 & 3 years & \$15{,}000 & \$75{,}000 & \$90{,}000 & \$90{,}000 & \$0 & \$0 \\
 & 4 years & \$15{,}000 & \$100{,}000 & \$115{,}000 & \$120{,}000 & \$5{,}000 & \$1{,}250 \\
 & 5 years & \$15{,}000 & \$125{,}000 & \$140{,}000 & \$150{,}000 & \$10{,}000 & \$2{,}000 \\
\bottomrule
\end{tabular}
\end{frame}

% ============================================================
\section{Phenomena Explained by Fixed Costs}
% ============================================================

\begin{frame}{Some Phenomena Explained by Fixed Costs}
\transfade
\begin{itemize}[<+->]
  \item Overtime
  \item Temporary help agencies
  \item Layoffs
  \item Segmentation of labour markets
  \item Resistance to work-sharing
\end{itemize}
\end{frame}

\begin{frame}{Figure 6.3: Nonrecurring Fixed Employment Costs and Changes in Labour Demand — \textit{Placeholder}}
\transfade
\placeholderbox{Figure 6.3 — Labour-demand shifts with nonrecurring fixed employment costs.}
\end{frame}

% ============================================================
\section{Worksharing and Job Creation}
% ============================================================

\begin{frame}{Worksharing and Job Creation}
\transfade
\begin{itemize}[<+->]
  \item Overtime Restrictions and Job Creation
  \item Nonwage Benefit and Part-time work
  \item Short-time Work and the Great Recession of 2008--2009
  \item Subsidizing Employment Sharing
  \item Workers and the Demand for Worksharing
\end{itemize}
\end{frame}

% ============================================================
\section{Chapter Summary}
% ============================================================

\begin{frame}{Summary}
\transfade
\begin{itemize}[<+->]
  \item Labour compensation includes substantial \textbf{nonwage} benefits.
  \item \textbf{Quasi-fixed} costs matter for employment vs.\ hours choices.
  \item These costs create a \textbf{wedge} between VMP and the wage.
  \item Implications for overtime, layoffs, segmentation, and worksharing policies.
\end{itemize}
\end{frame}

\begin{frame}
\centering
\Huge{\textbf{Thank You!}} \\

\end{frame}

\end{document}
